% Additive audit proof. No frozen provider is edited or compiled.
\section{Complete root--Gamma bounds for the two attained-tail monic fibres}
\label{sec:atgpa-audit}

Throughout this section retain the original parameters
\[
 0<\delta<\tfrac12,\qquad \gamma>2,\qquad
 j=4l+1\ge257,\qquad s\in\{1,j\},
\]
and the original fixed conductor order $v$, with
$0\le v\le\Delta$. Put
\[
\begin{gathered}
 q=(j+1)^2,\quad q'=(j-7)^2,\quad
 \Delta=16j-48,\quad m=\Delta-v,\quad
 c_+=j/2,\quad c_-=j/2-4,\\
 D=\sqrt{\delta^2+\gamma^2},\quad A=jD,\quad
 \epsilon=2^{-32},\quad T=\epsilon q,\quad
 A\le T/2,\quad 0\le r\le q+32j.
\end{gathered}\tag{ATG1}
\]
The letter $A$ here is a real root-radius bound, not the original
conductor operator. No parameter or root is changed by this notation.
Define
\[
 \ell_s=(s-1)/4,\quad M_s=(2\pi)^{s/2},\quad
 \beta_s=\frac{M_s}{\sqrt2\,\Gamma(s/2)},\quad
 \vartheta=\pi/2,\quad h=1/q.
\tag{ATG2}
\]
The integers used as monic degrees are $b_+=r$ and $b_-=m+r$;
the two root counts are $n_+=q$ and $n_-=q'$; and
$Q_\pm=n_\pm+\ell_s$.

Here are all finite constants in the comparison:
\[
\begin{gathered}
 C_L=\frac\pi{\Gamma(3/4)^2},\qquad
 C_U=\frac{\Gamma(1/4)^2}{2\pi},\qquad
 K_q=\frac{\sqrt{2\pi}}{C_L}\frac{(2q+1)^2}{q}e^{-2q},\\
 E_j=\frac{q A^2}{T^2-A^2}+\log(1+K_q),\\
 S_\ell=\sum_{a=0}^{\ell-1}(2a+\tfrac12)^2
       =\frac{\ell(16\ell^2-12\ell-1)}{12},\\
 \rho_0=0,\qquad
 \rho_\ell=\frac{S_\ell}{T^2}+\log(1+K_q)\quad(\ell>0),\\
 L_h=C_L(\sin h)^{3/2},\qquad U_h=C_U(\sin h)^{-1/2},\\
 L_{b,Q}=b+\sqrt{b(b+2Q)},\qquad
 K_{b,Q}=2Q+1+2L_{b,Q},\\
 b^-_{n,b,s}=-E_j+\log L_h
              -K_{b,n+\ell_s}\log(1+h/\vartheta),\\
 b^+_{n,b,s}=E_j+\rho_{\ell_s}+\log U_h
              +K_{b,n+\ell_s}\log\frac1{1-h/\vartheta}.
\end{gathered}\tag{ATG3}
\]
The $E_j$ in both signs is exactly the common GPA root error, with
the original symbol $k$ replaced by its attained-tail index $j$.
In particular, the lower-root application still uses the unchanged
common outer scale $q$, cutoff $T$, and rate displacement $h$.

We prove below, for $\zeta\in\{+,-\}$, every $b\le2q$, and every
$f\in\mathbb C[S]$ of degree at most $b$, the whole-space inequalities
\[
\begin{split}
 \beta_s e^{b^-_{n_\zeta,b,s}}
  \int_{\mathbb R}|f(c_\zeta+iy)|^2
           |y|^{2(n_\zeta+\ell_s)}e^{-\vartheta|y|}\,dy
 &\le\int_{\mathbb R}|\chi_{j_\zeta}(c_\zeta+iy)
                 f(c_\zeta+iy)|^2\,dm_s(y)\\
 &\le\beta_s e^{b^+_{n_\zeta,b,s}}
  \int_{\mathbb R}|f(c_\zeta+iy)|^2
           |y|^{2(n_\zeta+\ell_s)}e^{-\vartheta|y|}\,dy,
\end{split}\tag{ATG4}
\]
where $j_+=j$, $j_-=j-8$, and the original source is
\[
 dm_s(y)=M_s\frac{2^{s/2}}{4\pi\Gamma(s/2)}
                 |\Gamma(s/4+iy/2)|^2\,dy.
\tag{ATG5}
\]
No restriction on the zeros, parity, coefficients, or extremality of
$f$ occurs in (ATG4).

\subsection{The exact original root and coefficient maps}

The complete original root polynomials, including every multiplicity,
are
\[
\begin{split}
 P_j(y)&=\prod_{a,t=0}^{j}
       [y-(2t-j)\gamma+i(2a-j)\delta],\\
 P_{j-8}(y)&=\prod_{a,t=0}^{j-8}
       [y-(2t-j+8)\gamma+i(2a-j+8)\delta].
\end{split}\tag{ATG6}
\]
Their exact original-coordinate identities are
$\chi_j(c_++iy)=i^qP_j(y)$ and
$\chi_{j-8}(c_-+iy)=i^{q'}P_{j-8}(y)$.
The respective involutions $(a,t)\mapsto(j-a,j-t)$ and
$(a,t)\mapsto(j-8-a,j-8-t)$ pair every root with its negative.
All roots have modulus at most $A$, and both degrees are even.
For fixed $b$, the map
\[
 (\Pi_{c,b}f)(y)=i^{-b}f(c+iy),\qquad
 \Pi_{c,b}^{-1}p(S)=i^bp((S-c)/i)
\tag{ATG7}
\]
is an exact linear isomorphism of the degree-at-most-$b$ coefficient
spaces. Its matrix entry in row $a$, column $t$ is
$\binom ta c^{t-a}i^{a-b}$ for $a\le t$, and zero otherwise.
It maps the original monic degree-$b$ affine fibre bijectively onto
the monic degree-$b$ fibre in $y$. The factors $i^{-b}$ and $i^{n_\zeta}$
have absolute value one; they therefore disappear only after the
squared norms are formed, without changing any norm or minimum.

\subsection{Full source mass and the positive Gamma recurrence}

The defining Euler integral is the classical formula recorded by
R.~A. Askey and R. Roy \cite[eq.~(5.2.1)]{human:dlmf5}.
For $a>0$, two such Euler integrals, the substitution $t=e^u$ in their
beta integral, and Fourier inversion give
\[
 \frac{\Gamma(a+ix)\Gamma(a-ix)}{\Gamma(2a)}
 =\int_{\mathbb R}e^{ixu}(2\cosh(u/2))^{-2a}\,du,
 \qquad
 \int_{\mathbb R}|\Gamma(a+ix)|^2dx
 =2\pi\Gamma(2a)2^{-2a}.
\tag{ATG8}
\]
To justify the inversion absolutely, rotate Euler's integration ray
through $\operatorname{sgn}(x)\theta$, with $0<\theta<\pi/2$.
The large connecting arc is bounded by a polynomial times an
exponential with negative real exponent; the small arc tends to zero
because $a>0$. Taking the modulus on the rotated ray proves
\[
 |\Gamma(a+ix)|\le e^{-\theta|x|}\Gamma(a)(\cos\theta)^{-a}.
\tag{ATG9}
\]
Thus the integrable Fourier transform in (ATG8) may be inverted at
$u=0$. With $a=s/4$ and $y=2x$, (ATG8) proves
$\int dm_s=M_s$, including its full original mass.

Let
$d\sigma(y)=|\Gamma(1/4+iy/2)|^2dy/(2\pi)$.
Integration by parts in Euler's integral gives
$\Gamma(z+1)=z\Gamma(z)$, the classical functional equation
\cite[eq.~(5.5.1)]{human:dlmf5}. Iterating exactly $\ell_s$ times gives
\[
 dm_s(y)=\beta_s\prod_{a=0}^{\ell_s-1}
                (y^2+(2a+\tfrac12)^2)\,d\sigma(y).
\tag{ATG10}
\]
Indeed each squared recurrence factor contributes
$(y^2+(2a+1/2)^2)/4$; since $s/2=2\ell_s+1/2$, the remaining
constant is precisely $\beta_s/(2\pi)$ before the Gamma square.
For $s=1$, $\ell_s=0$ and $\beta_1=1$, so
$\int d\sigma=\sqrt{2\pi}$.

The reflection identity used here is Euler's reflection formula,
as recorded by Askey and Roy \cite[eq.~(5.5.3)]{human:dlmf5}.
For completeness, the reflection product needed below follows
directly from the beta integral
$\Gamma(z)\Gamma(1-z)=\int_0^\infty t^{z-1}/(1+t)\,dt$,
$0<\Re z<1$. On the keyhole contour with argument in $(0,2\pi)$,
the radial integrals give the multiplier $1-e^{2\pi iz}$; the
residue at $-1$ is $e^{i\pi(z-1)}$; the two circular contributions
tend to zero as $O(r^{\Re z})$ and $O(R^{\Re z-1})$.
The residue theorem therefore evaluates that integral as
$\pi/\sin(\pi z)$. At $z=1/4+iy/2$ it yields
\[
 |\Gamma(1/4+iy/2)|^2|\Gamma(3/4-iy/2)|^2
       =\frac{2\pi^2}{\cosh(\pi y)}.
\tag{ATG11}
\]
The unrotated Euler integral bounds the second squared factor by
$\Gamma(3/4)^2$, and consequently
\[
 \frac{d\sigma}{dy}\ge C_Le^{-\pi|y|},\qquad
 \frac{dm_s}{dy}\ge\beta_s q^{2\ell_s}C_Le^{-2\pi q}
                 \quad(q\le y\le2q).
\tag{ATG12}
\]
Furthermore, without suppressing either original mass,
\[
 \frac{M_s}{\beta_s q^{2\ell_s}}
 =\sqrt2\,\Gamma(2\ell_s+1/2)q^{-2\ell_s}
 \le\sqrt{2\pi}.
\tag{ATG13}
\]
To see the inequality, expand the Gamma value as
$\sqrt\pi\prod_{a=0}^{2\ell_s-1}(a+1/2)$ and bound each factor
by $s/2\le q$.

\subsection{A bound on every polynomial direction, including the inner region}

For every $p\in\mathbb C[y]$ of degree at most $b$ and
$0<T\le q/2$,
\[
 \sup_{|y|\le T}|p(y)|^2
 \le\frac{(b+1)^2 10^{2b}}q\int_q^{2q}|p(u)|^2du.
\tag{ATG14}
\]
Here is a coefficient-space proof. The Rodrigues polynomials
$L_a$ are the classical Legendre polynomials; the formula and norm
are recorded by T.~H. Koornwinder, R. Wong, R. Koekoek,
R.~F. Swarttouw and W.~P. Reinhardt
\cite[eq.~(18.5.5), Table~18.5.1, and the Legendre and
shifted-Legendre rows of Table~18.3.1]{human:dlmf18}.
In precisely those conventions, the polynomials
$L_a(x)=(2^aa!)^{-1}(d/dx)^a(x^2-1)^a$ have squared norm
$2/(2a+1)$ on $[-1,1]$ and are orthogonal to lower degrees,
by $a$ integrations by parts. The boundary derivatives of orders
less than $a$ vanish. Their norm follows by pairing the leading
coefficient $(2a)!/(2^a(a!)^2)$ against $x^a$ and using
$\int_0^1 t^a(1-t)^a dt=(a!)^2/(2a+1)!$, itself obtained by
repeated integration by parts. Leibniz's rule gives
$L_a(x)=2^{-a}\sum_{t=0}^a\binom at^2(x-1)^{a-t}(x+1)^t$.
For $|x|\le4$, its modulus is at most
$(5/2)^a\sum_t\binom at^2\le10^a$.
Expand $p(qz)$ in the orthonormal basis
$\sqrt{2a+1}L_a(2z-3)$ on $1\le z\le2$.
For $|z|\le1/2$, Cauchy--Schwarz and
$\sum_{a=0}^b(2a+1)=(b+1)^2$ give (ATG14), after changing
variables back to $y$.

Let $P$ be either polynomial (ATG6), with degree $n\in\{q,q'\}$.
For $|y|\ge T$, pairing roots $\omega,-\omega$ gives
\[
 \left|\log\frac{|P(y)|^2}{|y|^{2n}}\right|
 \le-n\log(1-A^2/T^2)\le\frac{nA^2}{T^2-A^2}.
\tag{ATG15}
\]
This follows from $1-u\le|1-z|\le1+u$ for $|z|\le u<1$,
and $-\log(1-u)\le u/(1-u)$.
For all real $y$, the same pairing gives
$|P(y)|^2\le(y^2+A^2)^n$.
Set $B=\int_q^{2q}|y|^{2n}|p(y)|^2dm_s(y)$.
By (ATG12)--(ATG14), each of
\[
 \int_{|y|<T}|P(y)p(y)|^2dm_s(y),\qquad
 \int_{|y|<T}|y|^{2n}|p(y)|^2dm_s(y)
\]
is at most $\kappa B$, where
\[
 \kappa=\frac{M_s}{\beta_s q^{2\ell_s}C_L}
  \frac{(b+1)^2}{q}e^{2\pi q}10^{2b}
             \left(\frac{T^2+A^2}{q^2}\right)^n.
\tag{ATG16}
\]
The bounds $n\ge q/2$, $b\le2q$, $A\le T/2$ imply
\[
 \kappa\le\frac{\sqrt{2\pi}}{C_L}\frac{(2q+1)^2}{q}
 \exp\{q(2\pi+4\log10+\tfrac12\log(5/4)+\log\epsilon)\}
 \le K_q.
\tag{ATG17}
\]
For an exact elementary check of the last exponent,
$\pi<4$, $\log10<4\log2$, $\log(5/4)<1/4$, and
$\log2>2/3$ give
$2\pi+4\log10+\tfrac12\log(5/4)-32\log2
 <8-16\log2+1/8<-61/24<-2$.
The inequalities for $\pi$ and $\log2$ follow respectively from
$\pi/4=\int_0^1(1+x^2)^{-1}dx$ and
$\log2=2\int_0^{1/3}(1-t^2)^{-1}dt$.

Let $I_{\rm out}$ be the reference integral with $|y|^{2n}$
outside $(-T,T)$. Since $B\le I_{\rm out}$, both inner integrals
are bounded by $K_q I_{\rm out}$. The full reference norm lies
between $I_{\rm out}$ and $(1+K_q)I_{\rm out}$.
The full actual norm lies between
$e^{-nA^2/(T^2-A^2)}I_{\rm out}$ and
$e^{nA^2/(T^2-A^2)}(1+K_q)I_{\rm out}$.
Thus for every coefficient direction, including the zero polynomial,
\[
 e^{-E_j}\int|p|^2|y|^{2n}dm_s
 \le\int|pP|^2dm_s
 \le e^{E_j}\int|p|^2|y|^{2n}dm_s.
\tag{ATG18}
\]

We also prove the full-form recurrence comparison. Put $Q=n+\ell_s$.
Pointwise,
$|y|^{2n}\prod_a(y^2+(2a+1/2)^2)\ge |y|^{2Q}$.
For $|y|\ge T$ its logarithmic ratio to $|y|^{2Q}$ is at most
$S_{\ell_s}/T^2$. On $|y|<T$ its numerator is at most
$(T^2+(j/2)^2)^Q$.
Apply (ATG14), the $s=1$ case of (ATG12), and
$\int d\sigma=\sqrt{2\pi}$. Since $j/2\le A\le T/2$,
$Q\ge q/2$, and the base $(5/4)\epsilon^2$ is less than one,
the calculation (ATG16)--(ATG17) bounds the entire inner numerator
by $K_q$ times the reference integral on $[q,2q]$.
The full upper ratio is therefore at most
$e^{S_{\ell_s}/T^2}+K_q\le e^{\rho_{\ell_s}}$.
At $\ell_s=0$ the comparison is the exact identity. In all cases,
\[
 \beta_s\int|p|^2|y|^{2Q}d\sigma
 \le\int|p|^2|y|^{2n}dm_s
 \le\beta_s e^{\rho_{\ell_s}}\int|p|^2|y|^{2Q}d\sigma.
\tag{ATG19}
\]
In particular, no inner interval has been discarded in either
comparison.

\subsection{Global density bounds and their exact rate-change map}

The classical inputs are again Euler's integral and Euler's
reflection identity \cite[eqs.~(5.2.1), (5.5.3)]{human:dlmf5};
the complete inequalities following from them are proved here.
Apply (ATG9) with $\theta=\pi/2-h$, first at $a=1/4$.
After squaring and dividing by $2\pi$, it gives
$d\sigma/dy\le U_he^{-(\vartheta-h)|y|}$.
At $a=3/4$, (ATG9) bounds the second Gamma square in (ATG11)
by $\Gamma(3/4)^2(\sin h)^{-3/2}e^{-(\vartheta-h)|y|}$.
Dividing (ATG11) by that bound, then using
$\cosh(\pi y)\le e^{\pi|y|}$, proves on the whole real line
\[
 L_he^{-(\vartheta+h)|y|}
 \le\frac{d\sigma}{dy}(y)
 \le U_he^{-(\vartheta-h)|y|}.
\tag{ATG20}
\]

Write $\|p\|_{Q,\kappa}^2=\int|p(y)|^2|y|^{2Q}
e^{-\kappa|y|}dy$ and $\mathsf A p=yp'(y)$ on polynomials of
degree at most $b$. We claim
$\|\mathsf A\|\le L_{b,Q}$ for each $\kappa>0$.
On the positive half-line, set $x=\kappa y$ and $\alpha=2Q$.
The classical Laguerre Rodrigues formula, series, and squared norms
are recorded by Koornwinder, Wong, Koekoek, Swarttouw and Reinhardt
\cite[Table~18.3.1; eq.~(18.5.5), Table~18.5.1;
eq.~(18.5.12)]{human:dlmf18}. In those exact conventions, the polynomials
\[
 \mathcal L_a^{(\alpha)}(x)
 =\frac{x^{-\alpha}e^x}{a!}
          \left(\frac d{dx}\right)^a(e^{-x}x^{a+\alpha})
 =\sum_{t=0}^a\frac{(-1)^t}{t!}
                         \binom{a+\alpha}{a-t}x^t
\tag{ATG21}
\]
are orthogonal to lower degrees and have squared norm
$\Gamma(a+\alpha+1)/a!$ in $x^\alpha e^{-x}dx$.
These assertions follow by $a$ integrations by parts; every boundary
term vanishes, as the power at zero before the final integration is
at least $\alpha+1$ and the exponential controls infinity.
The coefficient formula in (ATG21) gives
$x(\mathcal L_a^{(\alpha)})'
=a\mathcal L_a^{(\alpha)}-(a+\alpha)\mathcal L_{a-1}^{(\alpha)}$.
In the orthonormal basis, this is the sum of a diagonal matrix of
norm $b$ and a one-step shift of norm $\sqrt{b(b+2Q)}$.
The same estimate applies to $p(-y)$ on the negative half-line;
adding the two squared estimates proves the claim on the full line.
The pullback $x=\kappa y$ preserves the operator and contributes
only the common positive factor $\kappa^{-2Q-1}$ to squared norms.

The exact dilation map is $p(y)\mapsto p(e^ty)=\exp(t\mathsf A)p$,
with diagonal monomial matrix $\operatorname{diag}(1,e^t,\ldots,e^{bt})$.
The finite matrix exponential bound, followed by the same bound for
its inverse, gives
$e^{-L_{b,Q}|t|}\|p\|_{Q,\vartheta}
\le\|p(e^t\cdot)\|_{Q,\vartheta}
\le e^{L_{b,Q}|t|}\|p\|_{Q,\vartheta}$.
With $u=\kappa/\vartheta$, the exact substitution is
\[
 \|p\|_{Q,\kappa}^2
 =u^{-(2Q+1)}\|p(u^{-1}\cdot)\|_{Q,\vartheta}^2.
\tag{ATG22}
\]
Retaining that scalar first and then bounding it gives
\[
 e^{-K_{b,Q}|\log(\kappa/\vartheta)|}\|p\|_{Q,\vartheta}^2
 \le\|p\|_{Q,\kappa}^2
 \le e^{K_{b,Q}|\log(\kappa/\vartheta)|}\|p\|_{Q,\vartheta}^2.
\tag{ATG23}
\]
Combining (ATG18)--(ATG20) and (ATG23), at rates
$\vartheta+h$ and $\vartheta-h$, proves (ATG4) with precisely
the constants (ATG3).

\subsection{Domain, explicit error control, and monic minima}

All uses of the preceding inequalities are in their proved domains.
Indeed
$2q'-q=j^2-30j+97>0$ for $j\ge257$, so $q'\ge q/2$;
$Q_\pm\le q+(j-1)/4\le2q$; and
\[
 r\le q+32j\le2q,\qquad
 0\le m+r\le q+48j-48\le2q.
\tag{ATG24}
\]
For the last two upper inequalities, subtract the left sides from
$2q$: the resulting lower bounds are respectively
$j^2-30j+1$ and $j^2-46j+49$, both positive for $j\ge257$.
Also $j/2\le A$ because $\gamma>2$, and $0<h<\vartheta$.
The nonnegativity of $m+r$ uses the original $m=\Delta-v\ge0$;
it is not inferred merely from the displayed numerical domain.

Set $B_q=\log(1+K_q)$ and
$G_q=\log q+\log(\pi/2)$. Then the following fully explicit
bounds apply simultaneously to both signs, both source orders,
both root counts, and all degrees $b\le2q$:
\[
\begin{split}
 |b^-_{n,b,s}|&\le
 \frac{4D^2}{3\epsilon^2}+B_q+|\log C_L|
       +\tfrac32G_q+\frac{17}{\vartheta-1},\\
 |b^+_{n,b,s}|&\le
 \frac{4D^2}{3\epsilon^2}+2B_q+
       \frac1{16\epsilon^2j}+|\log C_U|
       +\tfrac12G_q+\frac{17}{\vartheta-1}.
\end{split}\tag{ATG25}
\]
To verify these, $A^2\le T^2/4$ and $j^2/q\le1$ imply
$qA^2/(T^2-A^2)\le4D^2/(3\epsilon^2)$.
Further,
$S_{\ell_s}\le\ell_s j^2/4\le j^3/16$ and $q\ge j^2$.
Since $b,Q\le2q$,
$K_{b,Q}\le(8+4\sqrt3)q+1\le16q+1$.
Both absolute rate logarithms are at most
$h/(\vartheta-h)$; their products with $K_{b,Q}$ are at most
$(16q+1)/(\vartheta q-1)\le17/(\vartheta-1)$.
Finally $2h/\pi\le\sin h\le h$, so
$|\log\sin h|\le G_q$.
The exact $K_q$ is bounded, since
$K_q\le9\sqrt{2\pi}C_L^{-1}q e^{-2q}
\le9\sqrt{2\pi}C_L^{-1}e^{-2}$ for $q\ge1$.
This proves the uniform $O_{\delta,\gamma}(\log q)$ claim without
absorbing $\beta_s$ into the error.

The two monic minima are defined with their original centres by
\[
 \nu^+_{r,s}=\min_{f\text{ monic},\,\deg f=r}
       \int|\chi_j(c_++iy)f(c_++iy)|^2dm_s(y),\qquad
 \nu^-_{m+r,s}=\min_{f\text{ monic},\,\deg f=m+r}
       \int|\chi_{j-8}(c_-+iy)f(c_-+iy)|^2dm_s(y).
\tag{ATG26}
\]
Each minimum exists uniquely: all moments are finite by (ATG9),
and the Gram form is positive definite because a nonzero polynomial
cannot vanish on a real set of positive measure. Its restriction to
the finite-dimensional lower-coefficient vector is consequently
strictly positive definite and coercive.
Define
$\mathfrak n(Q,b)=\min_{p\text{ monic},\,\deg p=b}
\|p\|_{Q,\vartheta}^2$.
The exact affine-fibre bijections (ATG7) allow minimization of both
sides of (ATG4) over exactly these fibres. Therefore
\[
\boxed{
 \nu^+_{r,s}\le\beta_s e^{b^+_{q,r,s}}
                       \mathfrak n(q+\ell_s,r),\qquad
 \nu^-_{m+r,s}\ge\beta_s e^{b^-_{q',m+r,s}}
                       \mathfrak n(q'+\ell_s,m+r).
}
\tag{ATG27}
\]
Both opposite inequalities also hold with their corresponding
$b^-$ and $b^+$. The positive original mass cancels only upon
forming the matched same-source ratio:
\[
 \frac{\nu^+_{r,s}}{\nu^-_{m+r,s}}
 \le e^{b^+_{q,r,s}-b^-_{q',m+r,s}}
       \frac{\mathfrak n(q+\ell_s,r)}
            {\mathfrak n(q'+\ell_s,m+r)}.
\tag{ATG28}
\]
No relation Gram, coefficient direction, source order, or root
coefficient is changed in this passage. The proof holds on all of
(ATG1); restricting to
$\lfloor j^{3/2}/\log j\rfloor<r\le q+32j$ preserves every
constant and makes the claimed tail uniformity immediate.
